\section{Data and Reproducibility}
\label{sec:reproducibility}

\subsection{Artifact structure}
\begin{itemize}
  \item \path{exports/truth_decay_pilot/} --- RQ1--RQ4 tables, audits, figures.
  \item \path{exports/truth_pilot/p4_validation.md} --- attribution validation.
  \item \path{exports/rq5_agent_impact/} --- partial A/B runs, uptake, ABC analysis.
  \item \path{exports/rq5_agent_impact_c/} --- complete condition~C runs.
  \item \path{exports/paper_synthesis/late_binding_evidence_table.csv} --- claim-to-evidence map.
  \item \path{docs/LATE_BINDING_MODEL_v1.md} --- conceptual synthesis.
  \item \path{protocol/RQ5_AGENT_IMPACT_EXPERIMENT_v1.md} --- experiment protocol v1.1.
\end{itemize}

\subsection{Replication commands}
\todo{Insert exact Makefile targets from repository freeze commit.}
Example frozen invocations (no agent reruns required for observational track):
\begin{verbatim}
make truth-decay-rq1 truth-decay-rq2 ...
make truth-decay-rq5-redesign   # plan only
\end{verbatim}

\subsection{Experiment freeze statement}
As of paper-writing mode, \textbf{no further agent experiments are scheduled}.
RQ5 A/B partial (128 runs) and C complete (105 runs) are frozen.

\subsection{Evidence table maintenance}
Each row in \path{late_binding_evidence_table.csv} includes a \texttt{limitation} column.
Authors must not cite a row without its stated limitation in the narrative.
